% ch3_answers.tex -- Quick check answers for Module 5, Chapter 3.
\begin{chapteranswers}{3}{Circular Motion: Angles, Energy, Work and Torque}
  % ---- 3.1 Angles and angular velocity
  \item (a) $\ang{60} \times \dfrac{\pi}{180} = \dfrac{\pi}{3} = \qty{1.047}{rad}$, so
        \final{\qty{1.05}{rad}}.
        (b) $2.5 \times 2\pi = 5\pi = \qty{15.71}{rad}$, so \final{\qty{15.7}{rad}}.
        (c) $3000 \times \dfrac{2\pi}{60} = \qty{314.2}{rad.s^{-1}}$, so
        \final{$\omega = \qty{314}{rad.s^{-1}}$}.
  \item (a) Both points turn through the same angle in the same time, so they have the
        \emph{same} $\omega$. Neither is larger.
        (b) $\omega = 45 \times \dfrac{2\pi}{60} = \qty{4.712}{rad.s^{-1}}$. Then
        $v_\mathrm{A} = \omega r = 4.712 \times 0.050 = \qty{0.2356}{m.s^{-1}}$ and
        $v_\mathrm{B} = 4.712 \times 0.15 = \qty{0.7069}{m.s^{-1}}$.
        Answers: \final{$v_\mathrm{A} = \qty{0.24}{m.s^{-1}}$} and
        \final{$v_\mathrm{B} = \qty{0.71}{m.s^{-1}}$}. B is 3 times as far out, so it is 3 times as fast.
  \item $v = 72 \div 3.6 = \qty{20}{m.s^{-1}}$. The wheel does not slip, so the rim moves at the
        car's speed: $\omega = \dfrac{v}{r} = \dfrac{20}{0.30} = \qty{66.67}{rad.s^{-1}}$.
        In rpm: $66.67 \times \dfrac{60}{2\pi} = 636.6$.
        Answers: \final{$\omega = \qty{67}{rad.s^{-1}}$}, about \final{640 rpm}.
  % ---- 3.2 Work and energy in circular motion
  \item The string is horizontal, so the tension is the whole centripetal force.
        (a) $F_\mathrm{T} = F_\mathrm{c} = \dfrac{mv^2}{r} = \dfrac{0.050 \times 5.0^2}{0.60} =
        \qty{2.083}{N}$, so the tension is \final{\qty{2.1}{N}}.
        (b) The tension points to the centre, at \ang{90} to every small step of the motion, so
        $W = Fs\cos\ang{90} = 0$: \final{$W = \qty{0}{J}$}, however many turns.
        (c) $K = \tfrac{1}{2}mv^2 = \tfrac{1}{2} \times 0.050 \times 5.0^2 = \qty{0.625}{J}$, so
        \final{$K = \qty{0.63}{J}$}, the same all the way round.
  \item Gravity pulls the satellite towards the Earth's centre. In a circular orbit this force is
        at \ang{90} to the satellite's velocity at every instant. So gravity has no part along the
        motion and does no work ($W = Fs\cos\ang{90} = 0$). By the work--energy theorem
        $\Delta K = 0$, so $K$ and the speed stay constant. The orbit radius does not change, so the
        potential energy $U$ is also constant. Therefore the total energy $E = K + U$ stays constant.
  \item (a) At the smallest speed, the bucket's push on the water is zero. So gravity alone gives
        the centripetal force: $mg = \dfrac{mv^2}{r}$. Cancel $m$ and multiply by $r$: $v^2 = gr$,
        so $v = \sqrt{gr} = \sqrt{9.8 \times 1.1} = \qty{3.283}{m.s^{-1}}$.
        Answer: \final{$v_\mathrm{top} = \qty{3.3}{m.s^{-1}}$}.
        (b) Energy from the top to the bottom (fall $= 2r$):
        $v_\mathrm{bottom}^2 = v_\mathrm{top}^2 + 4gr = 10.78 + 4 \times 9.8 \times 1.1 = \qty{53.9}{m^2.s^{-2}}$,
        so $v_\mathrm{bottom} = \qty{7.342}{m.s^{-1}}$. Answer: \final{$v_\mathrm{bottom} = \qty{7.3}{m.s^{-1}}$}.
  \item (a) No. The speed is constant, so $\Delta K = 0$.
        (b) The rider rises by the diameter, $2r = \qty{30}{m}$:
        $\Delta U = mg\Delta h = 50 \times 9.8 \times 30 = \qty{14700}{J}$, so
        \final{$\Delta U = \qty{1.5e4}{J}$} (a gain).
        (c) The seat's push on the rider, powered by the wheel's motor. It does
        $W_\mathrm{other} = \Delta K + \Delta U = 0 + \qty{1.5e4}{J}$ of work. Gravity does
        $-\qty{1.5e4}{J}$, so the net work is zero.
  % ---- 3.3 Torque
  \item (a) $\tau = rF\sin\theta = 0.35 \times 120 \times \sin\ang{90} = \qty{42}{N.m}$, so
        \final{$\tau = \qty{42}{N.m}$}.
        (b) $\tau = 0.35 \times 120 \times \sin\ang{70} = 42 \times 0.9397 = \qty{39.47}{N.m}$, so
        \final{$\tau = \qty{39}{N.m}$}.
  \item The force is straight down; $\theta$ is the angle between the crank and the force.
        (a) Crank horizontal: $\theta = \ang{90}$, so $\tau = 0.17 \times 600 \times 1 = \qty{102}{N.m}$:
        \final{$\tau = \qty{1.0e2}{N.m}$}.
        (b) \ang{45} below the horizontal: the crank is \ang{45} from the vertical force, so
        $\tau = 0.17 \times 600 \times \sin\ang{45} = 102 \times 0.7071 = \qty{72.12}{N.m}$:
        \final{$\tau = \qty{72}{N.m}$}.
        (c) Straight down: the force acts along the crank ($\theta = \ang{0}$), so
        \final{$\tau = \qty{0}{N.m}$}.
  \item The seesaw balances when the torques about the pivot are equal:
        $\tau_\mathrm{child} = \tau_\mathrm{adult}$, so
        $m_\mathrm{child}\, g\, r_\mathrm{child} = m_\mathrm{adult}\, g\, r_\mathrm{adult}$.
        The $g$ cancels: $30 \times 1.8 = 45 \times r_\mathrm{adult}$, so
        $r_\mathrm{adult} = 54 \div 45 = \qty{1.2}{m}$ (divide both sides by 45).
        Answer: \final{$r = \qty{1.2}{m}$} from the pivot. The adult is heavier, so the adult must
        sit closer to the pivot. (Each torque is $30 \times 9.8 \times 1.8 = \qty{529}{N.m}$.)
\end{chapteranswers}
